\documentclass{article}

\usepackage[a4paper, left=1.5cm, right=1.5cm, top=2cm, bottom=2cm]{geometry}
\usepackage{../../../../components/components}

\usepackage{fancyhdr}
\usepackage{amsmath}
\usepackage{amssymb}
\usepackage{hyperref}

% Opérateur PGCD propre
\DeclareMathOperator{\pgcd}{pgcd}

% Configuration des en-têtes et pieds de page
\pagestyle{fancy}
\fancyhf{} 

\fancyhead[L]{Préparation CAPES}
\fancyhead[C]{Arithmétique}
\fancyhead[R]{2026-2027}

\fancyfoot[L]{Ewen Rodrigues de Oliveira}
\fancyfoot[R]{\thepage}

\begin{document}
\docTitle{Chapitre 3 : Résolution dans $\mathbb{Z}$ d'équations diophantiennes linéaires}

\section{Cadre général}

\definition{
    Une équation diophantienne est une équation polynomiale à coefficients entiers dont les solutions recherchées sont des entiers.\\
    On s'intéresse ici aux équations diophantiennes linéaires à deux inconnues de la forme :
    \[
        (E) : ax + by = c
    \]
    avec $a, b, c \in \mathbb{Z}$ et l'inconnue $(x, y) \in \mathbb{Z}^2$.
}

\theorem{Théorème}{Existence de solutions entières}{true}{
    Soient $a, b, c \in \mathbb{Z}$.\\
    L'équation diophantienne linéaire $ax + by = c$ admet au moins une solution dans $\mathbb{Z}^2$ si et seulement si $\pgcd(a, b) \mid c$.
}

\noindent{\textbf{Preuve :}\\
Notons $d = \pgcd(a, b)$.
\begin{itemize}
    \item[$\Rightarrow$] Supposons que l'équation $ax + by = c$ admette une solution entière $(x_0, y_0)$. \\
    Par définition du PGCD, $d \mid a$ et $d \mid b$. Il existe donc deux entiers $a'$ et $b'$ tels que $a = da'$ et $b = db'$. En substituant dans l'équation, on obtient :
    \[
        da'x_0 + db'y_0 = c \iff d(a'x_0 + b'y_0) = c
    \]
    Comme $(a'x_0 + b'y_0) \in \mathbb{Z}$, on en déduit immédiatement que $d \mid c$.
    
    \item[$\Leftarrow$] Réciproquement, supposons que $d \mid c$. Il existe donc un entier $m$ tel que $c = md$. \\
    D'après le théorème de Bézout, il existe un couple d'entiers $(u, v) \in \mathbb{Z}^2$ tel que :
    \[
        au + bv = d
    \]
    En multipliant cette égalité par $m$, on obtient :
    \[
        a(mu) + b(mv) = md = c
    \]
    Le couple $(x_0, y_0) = (mu, mv) \in \mathbb{Z}^2$ est donc une solution particulière de l'équation.
\end{itemize}
\hfill$\square$
}

\section{Résolution d'une équation diophantienne linéaire à deux inconnues}
\subsection{Algorithme d'Euclide étendu et solution particulière}

\method{Recherche d'une solution particulière}{
    Pour déterminer une solution particulière de l'équation $ax + by = c$ lorsque $\pgcd(a,b) \mid c$ :
    \begin{enumerate}
        \item Établir l'algorithme d'Euclide pour déterminer $d = \pgcd(a, b)$.
        \item Remonter l'algorithme (méthode des substitutions successives) pour exprimer $d$ sous la forme d'une combinaison linéaire de $a$ et $b$ : $au + bv = d$.
        \item Multiplier l'égalité obtenue par le quotient $\frac{c}{d}$ afin d'isoler $c$ et d'identifier par identification un couple solution particulière $(x_0, y_0)$.
    \end{enumerate}
}

\example{
    Considérons l'équation $23x + 17y = 3$.
    \begin{enumerate}
        \item \textbf{Algorithme d'Euclide :}
        \[
            23 = 1 \cdot 17 + 6 \quad (\text{L}_1)
        \]
        \[
            17 = 2 \cdot 6 + 5 \quad (\text{L}_2)
        \]
        \[
            6 = 1 \cdot 5 + \textbf{1} \quad (\text{L}_3)
        \]
        \[
            5 = 5 \cdot 1 + 0
        \]
        Ainsi, $\pgcd(23, 17) = 1$. Comme $1 \mid 3$, l'équation admet des solutions entières.
        
        \item \textbf{Remontée de l'algorithme :} \\
        Via $(\text{L}_3)$ : $1 = 6 - 1 \cdot 5$ \\
        Via $(\text{L}_2)$ : $1 = 6 - 1 \cdot (17 - 2 \cdot 6) = 3 \cdot 6 - 1 \cdot 17$ \\
        Via $(\text{L}_1)$ : $1 = 3 \cdot (23 - 1 \cdot 17) - 1 \cdot 17 = 3 \cdot 23 - 4 \cdot 17$
        
        On obtient l'identité de Bézout : $23 \cdot (3) + 17 \cdot (-4) = 1$.
        
        \item \textbf{Obtention de la solution particulière :} \\
        En multipliant par $3$, on obtient :
        \[
            23 \cdot (9) + 17 \cdot (-12) = 3
        \]
        Le couple $(x_0, y_0) = (9, -12)$ est une solution particulière de notre équation.
    \end{enumerate}
}

\subsection{Résolution générale par Analyse-Synthèse}

\method{Détermination de l'ensemble des solutions}{
    La recherche de la forme générale des solutions repose rigoureusement sur un raisonnement par **analyse-synthèse** :
    \begin{enumerate}
        \item \textbf{Analyse :} On suppose qu'un couple $(x,y)$ est solution. Par soustraction membre à membre avec la solution particulière, on se ramène à une équation homogène. L'application du \textbf{théorème de Gauss} fournit les conditions nécessaires sur la forme de $x$ et $y$.
        \item \textbf{Synthèse :} On vérifie réciproquement que tout couple répondant à cette forme structurelle est effectivement solution (condition suffisante).
    \end{enumerate}
}

\example{
    Résolvons complètement dans $\mathbb{Z}^2$ l'équation $(E) : 23x + 17y = 3$, sachant que $(9, -12)$ est solution particulière.
    
    \noindent{\textbf{Analyse :}\\
    Soit $(x, y) \in \mathbb{Z}^2$ un couple solution de $(E)$. On a le système :
    \[
        \begin{cases}
            23x + 17y = 3 \\
            23 \cdot 9 + 17 \cdot (-12) = 3
        \end{cases}
    \]
    Par soustraction membre à membre, on obtient l'équation homogène associée :
    \[
        23(x - 9) + 17(y + 12) = 0 \iff 23(x - 9) = -17(y + 12) \quad (E_0)
    \]
    De l'équation $(E_0)$, on déduit que $17 \mid 23(x - 9)$. \\
    Or, $\pgcd(17, 23) = 1$. D'après le \textbf{théorème de Gauss}, $17 \mid (x - 9)$. \\
    Il existe donc un entier relatif $k \in \mathbb{Z}$ tel que :
    \[
        x - 9 = 17k \iff x = 17k + 9
    \]
    En substituant $x-9$ par $17k$ dans $(E_0)$, il vient :
    \[
        23 \cdot (17k) = -17(y + 12) \iff 23k = -(y + 12) \iff y = -23k - 12
    \]
    Ainsi, si $(x,y)$ est solution, il est nécessairement de la forme $(17k + 9, -23k - 12)$ avec $k \in \mathbb{Z}$.
    }

    \noindent{\textbf{Synthèse :}\\
    Soit $k \in \mathbb{Z}$. Considérons le couple $(x, y) = (17k + 9, -23k - 12)$. Injectons-le dans le membre de gauche de $(E)$ :
    \[
        23(17k + 9) + 17(-23k - 12) = 23 \cdot 17k + 207 - 17 \cdot 23k - 204 = 207 - 204 = 3
    \]
    Le couple est donc validé comme solution pour tout $k \in \mathbb{Z}$.
    }
    
    \noindent\textbf{Conclusion :} L'ensemble des solutions de l'équation dans $\mathbb{Z}^2$ est :
    \[
        \mathcal{S} = \left\{ (17k + 9 \ ; \ -23k - 12), \ k \in \mathbb{Z} \right\}
    \]
}

\newpage
\tableofcontents

\end{document}